\documentclass[tikz,border=12pt]{standalone}
\usepackage{amsmath}
\usepackage{kotex}
\usetikzlibrary{arrows.meta,positioning,shapes.geometric}
\definecolor{tensor}{RGB}{35,100,175}
\definecolor{model}{RGB}{128,75,160}
\definecolor{safe}{RGB}{30,130,95}
\definecolor{yield}{RGB}{198,68,54}
\begin{document}
\begin{tikzpicture}[>=Latex,font=\sffamily,line cap=round,line join=round]
  \node[font=\Large\bfseries] at (7.0,6.0) {다축 응력 상태를 일축 항복응력과 비교하는 방법};

  \node[font=\large\bfseries,tensor] at (2.1,5.15) {1. 응력텐서};
  \node[draw=tensor,very thick,rounded corners=3pt,fill=tensor!7,minimum width=4.0cm,minimum height=3.2cm,align=center] (T) at (2.1,3.05)
    {$\displaystyle[\boldsymbol\sigma]=\begin{bmatrix}
      \sigma_{11}&\sigma_{12}&\sigma_{13}\\
      \sigma_{12}&\sigma_{22}&\sigma_{23}\\
      \sigma_{13}&\sigma_{23}&\sigma_{33}
    \end{bmatrix}$\\[7pt]
    수직응력과 전단응력이\\함께 작용하는 상태};

  \node[font=\large\bfseries,model] at (7.0,5.15) {2. 항복모델};
  \node[draw=model,very thick,rounded corners=3pt,fill=model!7,minimum width=3.5cm,minimum height=1.15cm,align=center] (M1) at (7.0,3.85)
    {Tresca\\최대 주응력 차};
  \node[draw=model,very thick,rounded corners=3pt,fill=model!7,minimum width=3.5cm,minimum height=1.15cm,align=center] (M2) at (7.0,2.25)
    {von Mises\\편차응력의 크기};

  \node[font=\large\bfseries,yield] at (11.9,5.15) {3. 스칼라 비교};
  \node[draw=yield,very thick,rounded corners=3pt,fill=yield!6,minimum width=3.7cm,minimum height=3.2cm,align=center] (C) at (11.9,3.05)
    {$\sigma_{eq}(\boldsymbol\sigma)$ 계산\\[8pt]
     $\sigma_{eq}<\sigma^Y$\;: \textcolor{safe}{탄성 범위}\\[5pt]
     $\sigma_{eq}=\sigma^Y$\;: \textcolor{yield}{항복 경계}};

  % 하나의 응력텐서에서 두 항복모델 중 하나를 선택한다.
  \coordinate (split) at (4.65,3.05);
  \draw[very thick] (T.east)--(split);
  \fill (split) circle (1.8pt);
  \draw[->,very thick] (split)|-(M1.west);
  \draw[->,very thick] (split)|-(M2.west);

  % 각 모델이 계산한 등가응력을 같은 기준값과 비교한다.
  \draw[->,very thick] (M1.east)--(10.05,3.85);
  \draw[->,very thick] (M2.east)--(10.05,2.25);

  \node[align=center,text=gray!80!black,font=\small] at (7.0,0.55)
    {등가응력은 텐서 자체가 아니라, 선택한 항복모델이 다축 응력 상태에 부여한 하나의 비교값이다.\\같은 응력텐서라도 모델에 따라 $\sigma_{eq}$가 달라질 수 있다.};
\end{tikzpicture}
\end{document}
